% ======================================================================
\section{Introduction}
\IEEEPARstart{L}{oad} forecasting underpins almost every decision a system operator makes:
unit commitment, reserve procurement, market clearing, network switching and demand-response
dispatch all consume a forecast of demand over horizons from one hour to several weeks
\cite{hong2016tutorial,hong2020review}. Errors are expensive and asymmetric. On a system the size of Germany's a percentage point of
day-ahead error is several hundred megawatts, and errors concentrated at the daily peak are
costlier still, because that is where the marginal generator is most expensive and the network most
stressed.

The field has adopted deep sequence models wholesale. Recurrent and convolutional hybrids
built for load \cite{kong2019lstm,salam2021inception}, Transformer variants designed for
long-horizon forecasting \cite{zhou2021informer,wu2021autoformer,nie2023patchtst,%
liu2024itransformer,wang2024timexer}, deliberately simple linear models
\cite{zeng2023dlinear,li2023rlinear,xu2024fits}, multi-scale and all-MLP mixers
\cite{wang2024timemixer,chen2023tsmixer}, basis-expansion networks
\cite{oreshkin2020nbeats,challu2023nhits} and, most recently, pre-trained time-series
foundation models \cite{ansari2024chronos,das2024timesfm} now define the state of the art on
the standard electricity benchmarks. This literature inherits an assumption from generic
sequence modelling that is false for electricity demand: that the periodic structure of the
signal is unknown and must be learned from data, in coordinates relative to whatever window
the model happens to be looking at.

Electricity load is not a generic sequence. It is the aggregate of millions of human
schedules synchronised to a shared clock and a shared calendar. Its dominant spectral mass
sits at frequencies that are known exactly---one cycle per day, per week, per year---and the
calendar of the forecast horizon, including its public holidays, is known before the forecast
is made. So, in every control room, is a numerical weather forecast. Three consequences follow.

\emph{Phase.} A model that encodes time as an offset from the start of its lookback window has
no way to distinguish 08:00 on a Monday from 08:00 on a Sunday except by inferring it from the
data, and a small error in that inference compounds across the cycles a long forecast must
span. Frequency-domain methods \cite{zhou2022fedformer,xu2024fits} inherit the same problem
in a different guise: the discrete Fourier basis of a finite window is anchored to the window,
not to the calendar.

\emph{Adaptation.} The amplitude and shape of the daily profile drift with efficiency
programmes, electrification and photovoltaic self-consumption, and they move abruptly every
December, when a fortnight of holidays makes weekdays behave like Sundays. A seasonal
representation baked into network weights at training time cannot follow either kind of move
without continued training after deployment \cite{kim2025tta}.

\emph{Evaluation.} Progress is measured in mean squared error on z-scored data, on a fixed
lookback, with each method given whatever inputs its architecture happens to accept. On
national load that practice compares information sets rather than methods: a model that
consumes weather is not comparable with one that cannot, and neither is comparable with the
operator's forecast, which consumes both weather and calendar. It also ignores uncertainty,
even though reserve sizing is a decision about a distribution \cite{gneiting2007scoring}.

\smallskip\noindent\textbf{Contributions.} This paper makes three, each tied to an experiment
that could have falsified it.

\begin{enumerate}
\item \textbf{A calendar tier estimated per window in closed form.} A harmonic design matrix
indexed by wall-clock time and extended with weekend and public-holiday day-types has its
coefficients obtained per window and channel by a weighted ridge problem solved in closed
form inside the forward pass; a residual network models what the tier leaves. The tier is
independent of that network: the same tier improves both a patch Transformer and a
doubly-residual basis stack, by amounts that differ sharply between the two reference systems and, on the German
one, shrink with the horizon (Section~\ref{sec:ablation}). Two details of the first design,
locking the basis to the calendar rather than to the window and evaluating it at local civil
time rather than UTC, are measured rather than claimed: the first is worth about one per
cent, the second nothing at hourly resolution.

\item \textbf{Where per-window estimation pays over a learned calendar map.} A learned
calendar embedding and a closed-form per-window fit are statistically indistinguishable when
training data are plentiful and stationary. We test the three conditions under which the
solver should differ---scarce training data, transfer to an unseen system without
fine-tuning, and temporal drift without retraining---against a controlled competitor that
receives the same design matrix as learned tokens, and report where it does and does not
(Section~\ref{sec:mech}); a second control, the leading model fed through the same covariate
head as ours, separates the covariate pathway from the architecture.

\item \textbf{An information-set-controlled operational benchmark.} Eight national systems
(DE, FR, GB, ES, IT, NL, PL, BE), 2015--2020, with ERA5 weather, national holidays and each
operator's own published day-ahead forecast. Every model---\nv{n.models.learned} learned
forecasters spanning classical, linear, recurrent, convolutional, Transformer, basis-expansion
and covariate-native families, wrapped from their official implementations where those exist,
plus two zero-shot foundation models---receives the identical inputs, is tuned under one
budget, and is reported over three seeds with Model Confidence Sets; the headline models are
run under three information regimes (load only; plus calendar and holidays; plus weather
degraded to operational forecast skill). The same protocol yields calibrated intervals under adaptive conformal
inference and a like-for-like comparison against eight transmission operators.
\end{enumerate}

We also report what did not hold. Two hypotheses we held at the outset---that the closed-form
solve beats a learned calendar map on plentiful data, and that gradient-free adaptation alone
delivers robustness to regime shift---are not supported by the experiments and are reported
as negative results rather than dropped; the ablation that decides the model's own
composition is run to the standard the rest of the paper demands; and the proposed model is
not the most accurate point forecaster in its own benchmark. A covariate-native MLP is, for a
reason the regime experiments make specific: the two models receive the same weather forecast
and differ in the pathway by which it enters, and the paper says so in its first results
paragraph rather than its last.
